% Part 1: Algebraic foundations: vectors, tensors, forms and multivectors | Part: 1 | Depends on: none | Status: done | Original lines: 318-344
%!TEX root = ../main.tex
% =====================================================================
\part{Algebraic foundations: vectors, tensors, forms and multivectors}\label{part:algebra}
% =====================================================================

Everything in this part happens in a single finite-dimensional real vector space $V$ (later the tangent space
$T_pM$ at a point of a manifold). We begin with the two objects that live there, arrows and the machines that
measure them, and with the one operation that connects them: \emph{contraction}. From the freest possible way of
multiplying vectors -- the tensor algebra -- we then obtain, by imposing relations, the three algebras that will carry the rest of the
document:
\[
\underbrace{V,\ V^*}_{\S\ref{sec:vectors}}\ \to\
\underbrace{\otimes,\ T(V)}_{\S\ref{sec:tensors}}\ \to\
\underbrace{g,\ Q}_{\S\ref{sec:metric}}\ \to\
\left\{\begin{array}{l}
 T(V)\xrightarrow{\ v\otimes v=0\ }\textstyle\bigwedge V\quad(\S\ref{sec:ext})\\[3pt]
 T(V)\xrightarrow{\ v\otimes v=Q(v)\ }\Cl(V,Q)\quad(\S\ref{sec:clifford}\text{--}\ref{sec:duality})
\end{array}\right.
\]

\begin{intuition}[the big idea of Part I]
Start with the \emph{freest possible} way to multiply vectors (the tensor algebra, in which no relations hold between
products). Then \emph{decide} which relations you want. If you demand that a vector squared be zero, you get areas and
volumes with orientation (the exterior algebra: metric-free, the home of differential forms). If you demand that a vector
squared be its squared length, you get an algebra that simultaneously encodes lengths, angles, areas, reflections and
rotations (geometric algebra: it needs a metric). Everything else in this part is a consequence of that one choice.
\end{intuition}
